% Fibonacci-Pell nearest gaps: an all-exponent classification and
% quadratic-unit orbit rigidity
%
% Authors:
%   Dr. Denys Dutykh (Mathematics Department, Khalifa University of Science
%   and Technology, Abu Dhabi, UAE)
%   Prof. Laurent Vuillon (Univ. Savoie Mont Blanc, CNRS, LAMA, Chambery,
%   France)
%
\section{Introduction}
\label{sec:introduction}

How often can the square of a Fibonacci quadratic integer lie just below a
power of the fundamental Pell unit, with the small difference pointing in a
primitive negative-norm unit direction? This question combines an
Archimedean nearest-neighbour condition with an exact norm--content condition.
The former occurs with positive density; the latter will leave only two
rows.

Let

\begin{equation}
 \lambda=1+\sqrt2,
 \qquad
 D_{q,\sigma}=\sigma F_q+F_{q+1}\sqrt2,
 \qquad
 q\geqslant1,
 \qquad
 \sigma\in\{1,-1\}.
 \label{eq:intro-D}
\end{equation}

Suppose that a positive integer $n$ lies in the strict factor-two window

\begin{equation}
 D_{q,\sigma}^2<\lambda^n<2D_{q,\sigma}^2.
 \label{eq:intro-window}
\end{equation}

There is at most one such $n$, because the ratio of two distinct integral
powers of $\lambda$ is at least $\lambda>2$. Write

\begin{equation}
 \Delta_{q,\sigma}=\lambda^n-D_{q,\sigma}^2,
 \qquad
 g_{q,\sigma}=\cont{\Delta_{q,\sigma}},
 \label{eq:intro-gap}
\end{equation}

where

\[
 N(a+b\sqrt2)=a^2-2b^2
\]

is the quadratic norm and the coefficient content of
$a+b\sqrt2\in\Z[\sqrt2]$ is $\gcd(\abs a,\abs b)$. The target condition is

\begin{equation}
 \norm{\Delta_{q,\sigma}}=-g_{q,\sigma}^2.
 \label{eq:intro-target}
\end{equation}

After division by its largest rational coefficient factor, the gap is then
a unit of norm $-1$. The two surviving examples already display the two possible
orientations:

\begin{align}
 \lambda^2-(-F_2+F_3\sqrt2)^2
 &=-6+6\sqrt2=6\lambda^{-1},
 \label{eq:intro-example-minus}\\
 \lambda^6-(F_4+F_5\sqrt2)^2
 &=40+40\sqrt2=40\lambda.
 \label{eq:intro-example-plus}
\end{align}

Our first main theorem says that there are no others, even when the parity
of $n$ is not prescribed.

\begin{theorem}[All-exponent nearest-gap classification]
\label{thm:classification-intro}
For $q\geqslant1$, $\sigma\in\{1,-1\}$, and positive $n$, the two conditions~\labelcref{eq:intro-window,eq:intro-target} hold exactly
in the two rows

\begin{equation}
 \begin{array}{ccccc}
 \toprule
 q&\sigma&n&g_{q,\sigma}&\Delta_{q,\sigma}/g_{q,\sigma}\\
 \midrule
 2&-1&2&6&\lambda^{-1}\\
 4& 1&6&40&\lambda\\
 \bottomrule
 \end{array}
 \label{eq:two-solutions}
\end{equation}

In particular, no odd exponent is a target. Every hypothetical target obeys
the intrinsic selector

\begin{equation}
 n\text{ odd}\Longrightarrow3\mid q,
 \qquad
 n\text{ even}\Longrightarrow3\nmid q.
 \label{eq:intro-parity-selector}
\end{equation}
\end{theorem}

The boundary $q\geqslant1$ is sharp: at $q=0$ one has
$D_{0,\sigma}=\sqrt2$ and $\lambda-D_{0,\sigma}^2=\lambda^{-1}$.

The proof is an exact, finite certificate built on an unbounded argument.
The parity selector first divides the proof into two genuinely different
branches. For even $n$, algebraic reconstruction turns the coefficient
content into a primitive factor rectangle. This gives the missing height
separation $h<n<2q$ for the negative unit exponent $-h$. Two linear forms in
logarithms then give $h=O(\log q)$ and $q<10^{32}$; certified
rational-interval reductions yield $h<192$ and $q<90$, after which exact
recurrence arithmetic is exhaustive. For odd $n$, the selector forces
$3\mid q$. Its factorisation has no analogue of $h<n$, so the proof must be
reordered: a fixed-coefficient logarithmic form first bounds
$\min\{h,q\log_\lambda\varphi\}$, where $\varphi=(1+\sqrt5)/2$, and only then
can the moving-coefficient
form bound $q$. The same certified reductions give $h<48$ or $q<96$, and
the exact terminal list is empty.

The scarcity in \cref{thm:classification-intro} does not come from scarcity
of windows. \Cref{thm:candidate-density} proves that both even candidate
windows occur for natural density
$\log_\lambda\sqrt2-\eta=0.3300\ldots$, with $\eta$ the phase constant of
\cref{eq:window-phases}, and that the density statement
persists on every arithmetic progression. The exact norm--content condition,
not the Archimedean approximation, is the rigid step.

Our second main theorem packages the classification as a complete arithmetic
orbit result. Let

\begin{equation}
 \mathsf S(u,v)=(v,u+2v),
 \qquad
 \mathcal H(u,v)=u^2+v^2,
 \qquad
 z_{q,\sigma}=(F_{q+1},F_{q+1}+\sigma F_q),
 \label{eq:intro-orbit-data}
\end{equation}

and write $\lambda^r=C_r+P_r\sqrt2$. We call $d$ the \emph{anchor}: it is
the Pell index at time $t=0$ in the orbit equation below.

\begin{theorem}[Complete Fibonacci-core Pell-orbit classification]
\label{thm:orbit-intro}
For every $q\geqslant1$, $\sigma\in\{1,-1\}$, positive anchor $d$, and
$t\geqslant0$, one has

\begin{equation}
 \mathcal H(\mathsf S^tz_{q,\sigma})=P_{d+2t}
 \label{eq:intro-orbit-hit}
\end{equation}

if and only if one of the following alternatives holds:

\begin{equation}
\begin{array}{cccc}
 \toprule
 \sigma&q&d&t\\
 \midrule
 1&1&3&\text{arbitrary }t\geqslant0\\
 -1&1&1&0\\
 -1&2&3&0\\
 -1&4&5&0\\
 \bottomrule
\end{array}
\label{eq:intro-orbit-list}
\end{equation}

There is no hit at a positive even anchor. In particular, the
consecutive-Fibonacci core $z_{q,1}=(F_{q+1},F_{q+2})$ has no hit for
$q\geqslant2$.
\end{theorem}

The bridge behind this theorem is more general than the Fibonacci cores.
For a positive ordered integral seed $z=(u,v)$, the map

\begin{equation}
 A_z=(v-u)+u\sqrt2
 \label{eq:intro-seed-map}
\end{equation}

obeys $A_{\mathsf S^tz}=\lambda^tA_z$. Thus an integral linear dynamical
system is integrally conjugate to multiplication by a quadratic unit. An orbit hit
at anchor $d$ is equivalent, outside the canonical unit seeds, to the
anchored gap

\begin{equation}
 \lambda^{d-1}-A_z^2=g\lambda^{-(2t+1)},
 \label{eq:intro-anchored-gap}
\end{equation}

which automatically lies in the strict factor-two window. This equivalence
retains the anchor and time in both directions; it is not merely a numerical
analogy.

Several classical structures clarify why the proof works. The equations
$N(X+Y\sqrt2)=\pm1$ define the norm-one torus and the torsor of norm $-1$, so
coefficient content records the integral lattice multiple of a primitive
torsor point. The matrix of $\mathsf S$ lies in $\mathrm{GL}_2(\Z)$, and
its symmetric-square action preserves the Lorentzian identity
$2\mathcal H^2-\mathcal B^2-\mathcal Q^2=0$ used in the orbit bridge.
Two of the three involutions of the biquadratic field $\Q(\sqrt2,\sqrt5)$
carry the argument: one controls Pell norms, while the other separates the
dominant and conjugate Fibonacci terms in the logarithmic forms. Finally,
the local results combine $p$-adic rank lifting with an Archimedean irrational
rotation. These are operational connections: each supplies an invariant,
factorisation, height, or recurrence mechanism used below.

Interactions between the Fibonacci and Pell sequences have been studied from
several directions. Alekseyev determined the exact intersection of the two
sequences \cite{Alekseyev2011}, and Bravo, Herrera and Luca extended the
question to common values of their generalised analogues
\cite{BravoHerreraLuca2021}. Pomeo and Bravo measured how closely a Pell
number can approach a Fibonacci number \cite{PomeoBravo2024}, Luca and Zottor
classified the Fibonacci numbers occurring among the $Y$-coordinates of Pell
equations \cite{LucaZottor2023}, and Kafle, Luca, Montejano, Szalay and
Togb\'e treated $X$-coordinates that are products of two Fibonacci numbers
\cite{KafleEtAl2019}. Each of those questions asks for a value coincidence
between two sequences, or for an Archimedean bound on their difference. The
present question is different in kind. The
window~\labelcref{eq:intro-window} is only its Archimedean part, and the arithmetic lies
in the exact norm--content \cref{eq:intro-target} imposed on the gap
itself. That equation sees both coefficient coordinates of
$\Delta_{q,\sigma}$ at once, which is why a positive-density supply of
windows leaves exactly two rows.

The nearest-gap problem arose in an ongoing programme on the
Markoff--Frobenius uniqueness conjecture. A positive integral triple
$(a,b,c)$ with $a\leqslant b\leqslant c$ is a Markoff triple when

\begin{equation}
 a^2+b^2+c^2=3abc,
 \label{eq:markoff-equation}
\end{equation}

and the conjecture asks whether the largest entry $c$ determines the
unordered pair $\{a,b\}$. The problem has arithmetic, geometric,
word-combinatorial, and trace formulations; see Aigner \cite{Aigner2013}, the
integral-necklace model of Fisac \cite{Fisac2025}, and Reutenauer's account
of Christoffel words and Markoff numbers \cite{Reutenauer2026}. The present
paper neither proves nor disproves this conjecture. In particular, we do not
prove that a hypothetical Markoff collision must enter the Fibonacci-core
family in \cref{thm:orbit-intro}. The arithmetic theorem is complete on its
stated domain; the upstream reduction remains open.

The motivation does, however, fix the two sequences, and \cref{sec:two-arms}
proves it. Each branch of the Markoff tree carries a quadratic Perron unit,
the dominant root of $X^2-3m_0X+1$ for the Markoff number $m_0$ that the
branch keeps. What the orbit bridge of \cref{sec:zero-defect} consumes is
not that unit but an integral square root of it acting on two-square data,
and comparing traces shows that such a square root forces $3m_0-2$ to be a
perfect square. The Fibonacci and Pell arms, carrying $\varphi$ and
$\lambda$, are the only branches on which one exists. The pairing studied
here is therefore forced, not chosen.

The local theory also draws a sharp methodological boundary. A natural
plus-sign cross-gcd on genuine nearest windows contains arbitrary prescribed
powers of $7$ and $17$, and more generally simultaneous products from an
infinite class of split primes. A prime-rank parity criterion is exact, while
the split prime $241$ gives a sharp obstruction. These facts rule out any
proof based only on uniform boundedness, squarefreeness, or finitely many
compatible local rank clocks. They do not produce a norm--content target.

\Cref{fig:logic} displays the dependency structure and the one unproved
Markoff-facing arrow.

\begin{figure}[htbp]
 \centering
 \begin{tikzpicture}[
   node distance=6mm and 8mm,
   every node/.style={font=\small,align=center},
   proof/.style={draw=blue!55!black,rounded corners,fill=blue!4,
     inner sep=5pt,text width=35mm,minimum height=11mm},
   context/.style={draw=black!55,rounded corners,fill=black!3,
     inner sep=5pt,text width=35mm,minimum height=11mm},
   open/.style={draw=red!60!black,dashed,rounded corners,fill=red!2,
     inner sep=5pt,text width=35mm,minimum height=11mm},
   arrow/.style={-{Latex[length=2mm]},thick}
 ]
  \node[proof] (target) {norm--content target\\at any exponent};
  \node[context,left=of target] (density) {positive-density\\candidate windows};
  \node[context,right=of target] (clocks) {unbounded local\\prime-power clocks};
  \node[proof,below=of target] (selector) {intrinsic parity\\selector};
  \node[proof,below left=of selector] (even) {even branch:\\rectangle, heights,\\certified reduction};
  \node[proof,below right=of selector] (odd) {odd branch:\\reordered heights,\\certified reduction};
  \node[proof,below=18mm of selector] (classify) {two targets only;\\no odd target};
  \node[proof,below left=of classify] (semiconj) {seed/unit\\integral conjugacy};
  \node[proof,below right=of classify] (direct) {direct Pell hits\\at time zero};
  \node[proof,below=20mm of classify] (orbit) {complete Fibonacci-core\\orbit list};
  \draw[arrow] (target)--(selector);
  \draw[arrow] (selector)--(even);
  \draw[arrow] (selector)--(odd);
  \draw[arrow] (even)--(classify);
  \draw[arrow] (odd)--(classify);
  \draw[arrow] (classify)--(orbit);
  \draw[arrow] (semiconj)--(orbit);
  \draw[arrow] (direct)--(orbit);
  \node[open,below=of orbit] (markoff) {hypothetical\\Markoff collision};
  \draw[arrow,red!60!black,dashed] (markoff)--
    node[right,font=\scriptsize] {reduction open} (orbit);
 \end{tikzpicture}
 \caption{Logical architecture. Solid arrows are proof dependencies: the
 orbit list combines the independent seed/unit conjugacy, the all-exponent
 classification, and the direct time-zero Pell hits. Density and local clocks
 are contextual, not hypotheses. The dashed red arrow is the unproved
 Markoff-to-orbit reduction.}
 \label{fig:logic}
\end{figure}

The paper is organised as follows. \Cref{sec:system} develops the
quadratic-unit dictionary, the exact window densities, and the primitive
factor rectangle. \Cref{sec:leading,sec:height,sec:classification} prove the
even branch, and \cref{sec:odd-exponents} proves the odd exclusion.
\Cref{sec:zero-defect} gives the general seed/unit integral conjugacy and the
complete Fibonacci-core orbit theorem. \Cref{sec:local-clocks} develops the
local rank criterion and simultaneous prime-power recurrence. We close by
stating the exact Markoff-facing boundary and further directions.
\Cref{app:certificates} specifies the finite certificates and their coverage.
